\documentclass[10pt,letterpaper]{article}
\usepackage[lmargin=2.2in,rmargin=0.7in,tmargin=1.0in,bmargin=1.0in]{geometry}
\usepackage{notes}
\usepackage{listings}

\lstdefinestyle{lecturematlab}{
    language=Matlab,
    basicstyle=\ttfamily\footnotesize,
    keywordstyle=\color{blue},
    commentstyle=\color{gray},
    stringstyle=\color{red!60!black},
    breaklines=true,
    columns=fullflexible,
    keepspaces=true,
    showstringspaces=false,
    frame=single
}

\begin{document}

\begin{center}
\Large\textbf{DAY 41 -- From Differential Equations to Animations}\\
\normalsize Optional Lecture
\end{center}

%%%%%%%%%%%%%%%%%%%%%%%%%%%%%%%%%%%%%%%%%%%%%%%%%%%%%%%%%%%%%%%%%%%%%%%%
\section{Euler's Method}

\subsection{Why Numerical Solutions?}

We have learned several methods for solving differential
equations. However, many equations do not have solutions
that we can conveniently write using familiar functions.

For example, the nonlinear pendulum equation
\[
\theta''+\gamma\theta'+4\sin\theta=0
\]
is easy to write down but difficult to solve explicitly.

Instead of finding a formula, we can approximate
the solution at many closely spaced times.
These values let us plot the solution and animate
the motion.

%-------------------------------------------------------------
\subsection{Following the Tangent Line}

Consider the initial value problem
\[
x'=f(t,x),\qquad x(0)=x_0.
\]
Suppose we know the solution at time $t$.
For a small time step $h$, the tangent-line approximation
gives
\[
x(t+h)\approx x(t)+hx'(t).
\]
Substituting the differential equation,
\[
x(t+h)\approx x(t)+hf(t,x(t)).
\]

Let $t_n=nh$ and let $x_n$ approximate $x(t_n)$.
Then \textbf{Euler's method} is
\[
\boxed{
t_{n+1}=t_n+h,
\qquad
x_{n+1}=x_n+hf(t_n,x_n).
}
\]
In words:
\[
\text{new value}
=
\text{old value}
+
(\text{time step})(\text{current rate of change}).
\]

%-------------------------------------------------------------
\subsection{A Short Calculation}

\begin{example}[1]{
Use Euler's method with $h=0.1$ to approximate
$x(0.1)$ and $x(0.2)$:
\[
x'=-2x,\qquad x(0)=1.
\]
}

\paragraph{Step 1. Write the update formula.}
Since $f(t,x)=-2x$,
\[
x_{n+1}=x_n+0.1(-2x_n)=0.8x_n.
\]

\paragraph{Step 2. Take two steps.}
Starting from $x_0=1$,
\[
x_1=0.8(1)=0.8,
\qquad
x_2=0.8(0.8)=0.64.
\]
Thus,
\[
\boxed{x(0.1)\approx0.8,\qquad x(0.2)\approx0.64.}
\]

\paragraph{Step 3. Compare with the exact solution.}
The exact solution is $x(t)=e^{-2t}$, so
\[
x(0.2)=e^{-0.4}\approx0.6703.
\]
Euler's method gives an approximation, not the exact value.
\end{example}

%-------------------------------------------------------------
\subsection{Implementing the Method}

A loop repeats the same update many times:

\begin{lstlisting}[style=lecturematlab]
h = 0.1;
t = 0:h:2;

x = zeros(size(t));
x(1) = 1;

for n = 1:length(t)-1
    x(n+1) = x(n) + h*(-2*x(n));
end

figure
plot(t, x, 'o-', t, exp(-2*t), '--')
xlabel('Time')
ylabel('x(t)')
legend('Euler', 'Exact')
grid on
\end{lstlisting}

MATLAB indexing begins at $1$, so \texttt{x(1)}
stores the initial value $x_0$.

Try changing $h$ to $0.05$ and then $0.01$.
Over this fixed time interval, the approximation
improves as the step size decreases.

%%%%%%%%%%%%%%%%%%%%%%%%%%%%%%%%%%%%%%%%%%%%%%%%%%%%%%%%%%%%%%%%%%%%%%%%
\section{Simulating a Pendulum}

\subsection{Euler's Method for Systems}

For a system
\[
\mathbf{x}'=\mathbf{F}(t,\mathbf{x}),
\]
Euler's method has the same form:
\[
\boxed{
\mathbf{x}_{n+1}
=
\mathbf{x}_n+h\mathbf{F}(t_n,\mathbf{x}_n).
}
\]
We update every component using the values
at the current time.

For the pendulum, let $v=\theta'$. Then
\[
\begin{dcases}
\theta'=v,\\
v'=-\gamma v-4\sin\theta.
\end{dcases}
\]
The Euler updates are
\[
\boxed{
\begin{aligned}
\theta_{n+1}&=\theta_n+hv_n,\\
v_{n+1}&=v_n+h(-\gamma v_n-4\sin\theta_n).
\end{aligned}
}
\]
Both updates use $\theta_n$ and $v_n$.
Do not use the newly computed $\theta_{n+1}$
in the second equation.

%-------------------------------------------------------------
\subsection{One Step by Hand}

\begin{example}[2]{
For the pendulum
\[
\theta''+0.4\theta'+4\sin\theta=0,
\qquad
\theta(0)=\frac{\pi}{2},
\qquad
\theta'(0)=0,
\]
take one Euler step with $h=0.1$.
}

\paragraph{Step 1. Update the angle.}
\[
\theta_1=\theta_0+hv_0
=\frac{\pi}{2}+0.1(0)
=\frac{\pi}{2}.
\]

\paragraph{Step 2. Update the angular velocity.}
\[
v_1
=0+0.1\left(-0.4(0)-4\sin\frac{\pi}{2}\right)
=-0.4.
\]
Therefore,
\[
\boxed{
\theta(0.1)\approx\frac{\pi}{2},
\qquad
v(0.1)\approx-0.4.
}
\]
The negative velocity indicates that the bob
begins moving back toward the hanging position.

The unchanged angle after the first step is a feature
of the approximation: Euler uses the initial velocity,
which is zero.
\end{example}

%-------------------------------------------------------------
\subsection{Computing the Motion}

Use a smaller time step for the simulation:

\begin{lstlisting}[style=lecturematlab]
% Parameters and initial conditions
gamma = 0.4;
theta0 = pi/2;
v0 = 0;

% Time grid
h = 1/300;
T = 15;
t = 0:h:T;

% Allocate arrays
theta = zeros(size(t));
v = zeros(size(t));

theta(1) = theta0;
v(1) = v0;

% Euler updates
for n = 1:length(t)-1
    theta(n+1) = theta(n) + h*v(n);
    v(n+1) = v(n) ...
        + h*(-gamma*v(n) - 4*sin(theta(n)));
end
\end{lstlisting}

We now have three arrays:
\begin{itemize}
    \item \texttt{t}: the times;
    \item \texttt{theta}: the approximate angles;
    \item \texttt{v}: the approximate angular velocities.
\end{itemize}

%-------------------------------------------------------------
\subsection{Two Ways to Plot the Solution}

A time plot shows how the angle changes.
A phase portrait shows angle and angular velocity
together.

\begin{lstlisting}[style=lecturematlab]
figure

subplot(1,2,1)
plot(t, theta, 'LineWidth', 1.5)
xlabel('Time')
ylabel('\theta')
title('Angle versus time')
grid on

subplot(1,2,2)
plot(theta, v, 'LineWidth', 1.5)
xlabel('\theta')
ylabel('Angular velocity')
title('Phase trajectory')
grid on
\end{lstlisting}

For these parameters, the angle oscillates with
decreasing amplitude. The phase trajectory spirals
toward the origin.

%-------------------------------------------------------------
\subsection{Checking the Approximation}

A smooth-looking curve is not necessarily accurate.
Repeat the calculation with a smaller $h$ and compare.

Euler's method can introduce artificial energy.
For example, consider the undamped linear oscillator
\[
\theta'=v,\qquad v'=-4\theta.
\]
Its exact energy
\[
E=\frac12v^2+2\theta^2
\]
is constant. Euler's method instead gives
\[
E_{n+1}=(1+4h^2)E_n.
\]
Thus, the numerical energy grows even though
the true energy does not.

For more accurate simulations, MATLAB provides
solvers such as \texttt{ode45}. The following replaces
the Euler loop and supplies values at the same times:

\begin{lstlisting}[style=lecturematlab]
F = @(t,z) [z(2); ...
    -gamma*z(2) - 4*sin(z(1))];

options = odeset('RelTol',1e-8,'AbsTol',1e-10);

[t, Z] = ode45(F, t, [theta0; v0], options);

theta = Z(:,1);
v = Z(:,2);
\end{lstlisting}

The rest of our plotting and animation code works
with either numerical solution.

%%%%%%%%%%%%%%%%%%%%%%%%%%%%%%%%%%%%%%%%%%%%%%%%%%%%%%%%%%%%%%%%%%%%%%%%
\section{Animating the Pendulum}

\subsection{From Angle to Position}

A simulation gives us $\theta(t)$, but an animation
needs the position of the bob.

Place the pivot at $(0,0)$. For a rod of length $L$,
\[
\boxed{
X(t)=L\sin\theta(t),
\qquad
Y(t)=-L\cos\theta(t).
}
\]
At $\theta=0$, the bob is at $(0,-L)$,
directly below the pivot.

For the drawing, we can use $L=1$ as a visual scale:

\begin{lstlisting}[style=lecturematlab]
L = 1;
X = L*sin(theta);
Y = -L*cos(theta);
\end{lstlisting}

%-------------------------------------------------------------
\subsection{Preparing the Figure}

We will display the pendulum beside its phase trajectory.
A moving point on the phase portrait shows the
current angle and angular velocity.

\begin{lstlisting}[style=lecturematlab]
fig = figure('Color','w', ...
    'Position',[100 100 900 400]);

% Pendulum panel
ax1 = subplot(1,2,1);
hold on
axis equal
axis([-1.2 1.2 -1.2 1.2])
grid on
xlabel('X')
ylabel('Y')

plot(0,0,'k.','MarkerSize',20);

rod = plot([0 X(1)], [0 Y(1)], ...
    'k-', 'LineWidth',2);

bob = plot(X(1),Y(1),'ro', ...
    'MarkerFaceColor','r','MarkerSize',12);

clockLabel = title(ax1,'t = 0');

% Phase portrait panel
ax2 = subplot(1,2,2);
plot(theta,v,'Color',[0.7 0.7 0.7])
hold on
grid on
xlabel('\theta')
ylabel('Angular velocity')
title('Phase trajectory')

point = plot(theta(1),v(1),'ro', ...
    'MarkerFaceColor','r','MarkerSize',8);

axis tight
axis manual
\end{lstlisting}

The variables \texttt{rod}, \texttt{bob}, and
\texttt{point} refer to objects in the figure.
We can move these objects by changing their coordinates.

%-------------------------------------------------------------
\subsection{Updating the Picture}

Our simulation uses $300$ time steps per second.
Displaying every tenth point gives $30$ frames
per second.

\begin{lstlisting}[style=lecturematlab]
stride = 10;
fps = 1/(stride*h);

for n = 1:stride:length(t)
    set(rod, ...
        'XData',[0 X(n)], ...
        'YData',[0 Y(n)]);

    set(bob,'XData',X(n),'YData',Y(n));

    set(point, ...
        'XData',theta(n), ...
        'YData',v(n));

    set(clockLabel, ...
        'String',sprintf('t = %.2f',t(n)));

    drawnow
    pause(1/fps)
end
\end{lstlisting}

The loop moves the rod, the bob, and the phase-plane
point together.

The pause controls the preview speed approximately.
Rendering the figure also takes time, so the preview
may run slower than real time.

%%%%%%%%%%%%%%%%%%%%%%%%%%%%%%%%%%%%%%%%%%%%%%%%%%%%%%%%%%%%%%%%%%%%%%%%
\section{Saving and Exploring}

\subsection{Saving a Video}

To save the animation, create a video file before
the animation loop:

\begin{lstlisting}[style=lecturematlab]
movieFile = VideoWriter('pendulum.avi');
movieFile.FrameRate = fps;
open(movieFile);
\end{lstlisting}

Inside the loop, replace
\texttt{pause(1/fps)} with:

\begin{lstlisting}[style=lecturematlab]
writeVideo(movieFile,getframe(fig));
\end{lstlisting}

Keep \texttt{drawnow} immediately before capturing
each frame. After the loop, close the video:

\begin{lstlisting}[style=lecturematlab]
close(movieFile);
\end{lstlisting}

The file \texttt{pendulum.avi} is saved in MATLAB's
current folder. Keep the figure size fixed while
recording so all frames have the same dimensions.

%-------------------------------------------------------------
\subsection{Experiments}

Change one parameter at a time, rerun the simulation,
and predict the result before watching it.

\begin{enumerate}
    \item \textbf{Remove damping.}
    Set $\gamma=0$. Does the motion persist?
    Compare Euler's method with \texttt{ode45}.

    \item \textbf{Increase damping.}
    Compare $\gamma=0.4$, $\gamma=4$, and $\gamma=6$.
    Here, the small-angle natural frequency is
    $\omega=2$, so critical damping occurs at
    $\gamma=2\omega=4$.

    \item \textbf{Give the bob a large initial kick.}
    Set $\theta_0=0$ and $v_0=5$.
    Compare the motion with and without damping.

    \item \textbf{Test the numerical method.}
    Increase the Euler step size.
    Does the simulation still agree with your
    mathematical expectations?
\end{enumerate}

Numerical simulation connects three descriptions
of the same solution: the differential equation,
the phase trajectory, and the physical motion.

\end{document}